\section{训练误差为何总是乐观}
\label{sec:13-Machine-Learning/10-Learning-Theory/01}

一轮网格搜索跑完，屏幕上摊着四十组配置的训练曲线。你按验证指标排序，挑出最好的那一组，报告里写下训练错误率 \(0.4\%\)。模型上线两周，线上错误率稳定在 \(3.8\%\)。你依次排查了分布漂移、标注口径、特征泄漏，把每一条都堵住之后，数字仍然对不上。

剩下的那部分落差与工程无关。它来自一个更早的地方：报告里那个 \(0.4\%\) 是在四十组配置中挑出来的最小值，而挑选用的数据和计算它的数据是同一批。本节要说清这份乐观从哪里生出来，为什么它不随样本量消失，以及数学上有哪两条修补路径。

\subsection{你能算的那个数，和你想要的那个数，取的不是同一个期望}

设样本对 \((X, Y)\) 服从 \(\mathcal{X} \times \mathcal{Y}\) 上的分布 \(P\)，训练集

\[
S = \bigl((X_1, Y_1), \dots, (X_n, Y_n)\bigr)
\]

由 \(n\) 个独立同分布的观测组成。一个假设 \(h: \mathcal{X} \to \mathcal{Y}\) 配上损失 \(\ell(h(X), Y)\)，部署时你关心的量是真风险

\[
R(h) = \E\bigl[\ell(h(X), Y)\bigr],
\]

期望对一个来自同一分布、从未见过的新样本取。你能算出来的只有经验风险，也就是训练集上的平均损失：

\[
\widehat R(h) = \frac{1}{n} \sum_{i=1}^{n} \ell(h(X_i), Y_i).
\]

在 0-1 损失 \(\ell = \ind\{h(X) \ne Y\}\) 下，\(R(h) = P(h(X) \ne Y)\) 是真实错误概率，\(\widehat R(h)\) 是训练集上的错误比例。本节以 0-1 损失展开，结论的形式对任何取值有界的损失都成立。

泛化误差指的是差值 \(R(h) - \widehat R(h)\)。学习的目标是让 \(R\) 小，\(\widehat R\) 只是它的一个代理。代理可以贴地而目标依旧高悬，两者之间的距离就是本单元要研究和控制的对象。

\subsection{假设先定下来，训练误差就是一次诚实的抽样}

先取最干净的情形：假设 \(h\) 在你看到任何训练数据之前就已经固定。比方说你翻开一本教科书，抄下一条现成的分类规则。这时 \(\ell(h(X_i), Y_i)\) 是 \(n\) 个独立同分布、取值落在 \([0,1]\) 的随机变量，每个的期望恰好是 \(R(h)\)。于是 \(\widehat R(h)\) 成为 \(R(h)\) 的无偏估计——把``抽 \(n\) 个样本、算一次训练误差''这个实验重复很多遍，平均结果就是真风险。

\hyperref[sec:05-Probability/07-Transforms-and-Bounds/04]{``Chernoff 方法与 Hoeffding 集中界''}中的 Hoeffding 不等式给出这些独立有界变量的均值偏离期望的概率上限：

\[
P\bigl(\abs{\widehat R(h) - R(h)} \ge \varepsilon\bigr)
\le 2e^{-2n\varepsilon^2},
\]

等价地，以至少 \(1 - \delta\) 的概率

\[
\abs{\widehat R(h) - R(h)}
\le \sqrt{\frac{\log(2/\delta)}{2n}}.
\]

代入数字掂量一下力度。取 \(n = 10^4\)、\(\delta = 0.05\)，界宽约为 \(0.014\)。一个固定的分类器，在一万个独立样本上测一次，测得的错误率与真实错误率的差距以至少 \(95\%\) 的概率不超过一点四个百分点。分布无关、有限样本、指数收紧——``留出一个测试集来评估模型''之所以可信，全部数学根基就是这几行。

这份保证有两个不可让渡的前提。其一，\(h\) 不能以任何方式依赖这批评估样本。其二，评估样本与将来部署时遇到的样本来自同一分布。后者一旦被打破，\(R(h)\) 的定义本身就换了对象，任何界都无从谈起；前者则是本节余下篇幅的主题。

\subsection{一旦让数据参与挑选，无偏性立刻反向倾斜}

真实的训练流程几乎从不是``先定假设、再看数据''。经验风险最小化（ERM）返回的是

\[
\widehat h \in \argmin_{h \in \mathcal{H}} \widehat R(h),
\]

你最后写进报告的数字是 \(\widehat R(\widehat h)\)。它与上一小节的 \(\widehat R(h)\) 只差一顶帽子，统计性质却完全换了。同一批数据先被用来从 \(\mathcal{H}\) 里挑出 \(\widehat h\)，又被用来给 \(\widehat h\) 打分。Hoeffding 所依赖的独立性被训练流程本身破坏了：第 \(i\) 个样本一边参与决定谁胜出，一边参与给胜出者评分，两次使用共享同一份随机波动。

偏差的方向可以直接定出来。对任意固定的 \(h \in \mathcal{H}\)，ERM 的输出在训练集上不会输给它：

\[
\widehat R(\widehat h) \le \widehat R(h) \quad \text{对每个 } h \in \mathcal{H}.
\]

两边取期望，再用固定 \(h\) 时 \(\widehat R(h)\) 的无偏性，得到

\[
\begin{aligned}
\E\bigl[\widehat R(\widehat h)\bigr]
&\le \E\bigl[\widehat R(h)\bigr] = R(h) \quad \text{对每个 } h \in \mathcal{H},\\[2pt]
\E\bigl[\widehat R(\widehat h)\bigr]
&\le R(h^*) \le \E\bigl[R(\widehat h)\bigr],
\end{aligned}
\]

其中 \(h^* \in \argmin_{h \in \mathcal{H}} R(h)\) 是类内真风险最优者，末一步用了 \(R(\widehat h) \ge R(h^*)\) 这个恒成立的事实。训练误差的期望被压在类内最优真风险之下，而被选中假设的真风险期望被顶在它之上。夹在中间的 \(R(h^*)\) 把两者隔开，\(\widehat R(\widehat h)\) 因此不可能是 \(R(\widehat h)\) 的无偏估计。

这份向下的偏差与数据集的运气无关。ERM 的职责就是挑出训练集上表现最好的那个假设，而``训练集上表现最好''里面掺着这批特定样本的友好噪声。下面这张图把它画成了一句话。

\begin{center}
\begin{tikzpicture}[>=stealth]
% axes
\draw[->] (0,0) -- (6.4,0) node[below right] {\footnotesize 训练误差 \(\widehat R(h)\)};
\draw[->] (0,0) -- (0,5.4) node[above left] {\footnotesize 真风险 \(R(h)\)};
% diagonal
\draw[dashed,gray] (0,0) -- (5.2,5.2);
\node[gray,rotate=45,anchor=north] at (4.75,4.75) {\footnotesize \(R=\widehat R\)};
% band
\draw[gray!45] (0,0.9) -- (4.4,5.3);
\draw[gray!45] (0.9,0) -- (5.3,4.4);
% cloud of hypotheses
\foreach \x/\y in {1.1/1.6, 1.4/1.1, 1.8/2.4, 2.0/1.7, 2.3/2.9, 2.6/2.2,
                   2.9/3.4, 3.2/2.8, 3.5/3.9, 3.8/3.3, 4.1/4.4, 4.4/3.8,
                   4.7/4.9, 1.6/2.0, 3.0/2.4, 4.0/4.0}
  \fill[gray!70] (\x,\y) circle (0.055);
% the ERM pick
\fill[black] (0.55,1.5) circle (0.09);
\node[anchor=west] at (0.7,1.7) {\footnotesize \(\widehat h\)};
% gap arrows
\draw[<->] (0.25,0.55) -- (0.25,1.42);
\node[anchor=west] at (0.33,1.20) {\footnotesize 落差};
\draw[dotted] (0,1.5) -- (0.55,1.5);
\draw[dotted] (0.55,0) -- (0.55,0.55);
\end{tikzpicture}
\end{center}

\emph{每个灰点是一个候选假设，横纵坐标分别是它的训练误差与真风险；点云绕对角线上下抖动，而``取最左端''这个动作系统性地选中了抖到对角线上方的那一个。}

图里的落差不来自任何一个假设本身，它来自``取最左端''这个动作。点云的竖直抖动幅度由样本量决定，抖动本身随 \(n\) 增大而收窄；但只要选择和评估共用同一批数据，最左端的点就仍然偏在对角线上方。候选越多，点云在左侧铺得越远，被选中者的落差越大——哪怕 \(\mathcal{H}\) 里装的全是没有价值的假设，只要数量够多，其中总有一个碰巧在这批样本上显得漂亮。

\subsection{一百万枚均匀硬币能造出一枚假的偏币}

把机器学习的术语全部剥掉，这个效应可以用硬币重现。准备 \(m\) 枚均匀硬币，每枚代表一个候选假设；每枚抛 \(n\) 次，正面比例作为它的``训练表现''；从 \(m\) 枚中选出正面比例最低的那一枚，报告它的比例。单枚硬币的比例以半径 \(\sqrt{\log(2/\delta)/(2n)}\) 集中在 \(1/2\) 附近，这是 Hoeffding 对固定硬币的承诺。可是 \(m\) 枚里最极端的那一枚，其偏离量的量级是 \(\sqrt{\log m/(2n)}\)——被惩罚的对象从 \(\log(1/\delta)\) 换成了 \(\log m\)。代入 \(m = 10^6\)、\(n = 100\)，被选中的硬币正面比例大约 \(0.5 - 0.26 \approx 0.24\)。一枚完全均匀的硬币，仅仅因为是从一百万个候选里挑出来的最优者，就长出了偏币的外观。

硬币没有毛病，制造偏差的是候选数量。同一个道理在统计推断里早有精确写法：\hyperref[sec:11-Mathematical-Statistics/05-Hypothesis-Testing/04]{``多重比较改变错误的计量单位''}说明做一百次独立显著性检验几乎必然撞出假阳性，与从一百万个候选模型里几乎必然撞出虚高的训练表现是同一个机制。每多一次``看过数据再做决定''，随机性的同一份预算就被多花掉一次。

工程纪律把这条账翻译成一句话：测试集只能用一次。调超参数、筛特征、比较结构，凡是``看过数据再决定''的操作都要走验证集或重抽样（\hyperref[sec:11-Mathematical-Statistics/06-Resampling-and-Model-Selection/03]{``交叉验证评估的是完整训练流程''}）。评估数据每参与一次决定，就有一部分变成了训练数据，它对模型的打分随之带上前面那张图里的落差。

\subsection{断裂发生在 Hoeffding 的第一环，修复只有两种}

回看 Hoeffding 不等式的证明，第一步写的是``\(\ell(h(X_i), Y_i)\) 独立同分布、均值为 \(R(h)\)''。当 \(h = \widehat h\) 由样本决定时，第 \(i\) 项同时出现在``选出 \(\widehat h\)''与``评估 \(\widehat h\)''两个角色里，各项不再独立于 \(\widehat h\)，条件均值也不再等于 \(R(\widehat h)\)。后续的 Markov 步、Chernoff 步、指数衰减全部架在这一环上，第一环断了，整条链就失效。

第一种修复是把选择与评估在物理上分开：数据切成两份，一份专管挑假设，一份专管给挑出来的假设打分。评估时那个假设相对测试集是固定的，Hoeffding 的前提恢复。留出法和交叉验证走的都是这条路，牺牲掉的是挑选阶段可用的样本量。

第二种修复是把界加强到整个假设类上。仍然用同一批数据做选择，但要求

\[
\sup_{h \in \mathcal{H}} \abs{\widehat R(h) - R(h)}
\]

以高概率很小。这样一来，无论 ERM 挑中了谁，落差都已在控制之内。这条路线叫一致收敛，所需样本量取决于 \(\mathcal{H}\) 的复杂度而非被选中假设的复杂度，是下一节的主题。

两条路怎么选，取决于手上有什么。测试集充足、分布稳定时，物理拆分的界更紧、假设更弱、实现更简单；测试集稀缺或需要反复做模型选择时，一致收敛给出理论上的安全边际，只是常数通常很松，工程里仍以重抽样验证为主。往后的 margin 界、PAC-Bayes、稳定性与泛化的联系、双下降，都是对这条故事线的加细，而故事的第一句不会变：训练误差的乐观首先来自``同一批数据既用于选择又用于评估''，其次才轮到模型有多复杂。把这个顺序颠倒过来，就会在错的环节找解药。
